% ============================================================
% OP_GUT3b_hypercharge_recon_v2.tex  (PROPOSAL-ONLY; v2 per GPT
% review of v1). L1 + L2 computed together, compactly.
% GPT corrections incorporated: observed charge classes are INPUT
% (no "derives e_0 = e/3" framing); mandatory framing sentence
% included; the three nu^c branches separated explicitly in L2;
% N3 (coupling magnitude) kept entirely outside.
% NO preprint edits.
% ============================================================
\section*{OP-GUT3b reconnaissance v2: lattice interlock and the
$B{-}L$ flat direction (proposal-only)}

\subsection*{0. Scope and conventions}
Working frame: the approved N1--N4 separation.
Convention (N4, fixed once): $Q=T_3+Y/2$ with the recorded
assignments eq:hypercharges
($Y_Q=\tfrac13$, $Y_{u^c}=-\tfrac43$, $Y_{d^c}=\tfrac23$,
$Y_L=-1$, $Y_{e^c}=2$, and, where admitted, $Y_{\nu^c}$).
For each multiplet define the triality $t\in\{0,\pm1\}$
($\mathbf3\to+1$, $\bar{\mathbf3}\to-1$, singlet $\to0$) and the
duality $d\in\{0,1\}$ (weak doublet $\to1$, singlet $\to0$).
Inputs admitted as phenomenological data: the observed multiplet
list and charge classes.
Mandatory framing (applies to everything below):
\emph{this is a consistency requirement on completions, not a
derivation of the observed hypercharges.}

\subsection*{1. L1 computed: the centre-quotient interlock
congruence}
Consider the centre element
$g_*=(\omega\mathbf 1_3,\,-\mathbf 1_2,\,e^{i\pi Y})$,
$\omega=e^{2\pi i/3}$, acting on a multiplet with data $(t,d,Y)$
by the phase $\exp\!\bigl[2\pi i\bigl(\tfrac t3+\tfrac d2
+\tfrac Y2\bigr)\bigr]$.
$g_*$ acts trivially on the multiplet iff the
\emph{interlock congruence}
\[
\mathrm{(C\mbox{-}\Gamma)}\qquad
\frac{2t}{3}+d+Y\;\in\;2\mathbb Z
\qquad\Longleftrightarrow\qquad
3Y+3d+2t\equiv0 \pmod 6
\]
holds.
\emph{Verification on the recorded assignments} (all six
multiplets): $Q:\tfrac23+1+\tfrac13=2$;
$u^c:-\tfrac23+0-\tfrac43=-2$; $d^c:-\tfrac23+0+\tfrac23=0$;
$L:0+1-1=0$; $e^c:0+0+2=2$; $\nu^c(Y=0):0$.
All in $2\mathbb Z$; hence the observed one-generation content is
consistent with the maximal quotient $\Gamma=\mathbb Z_6$
generated by $g_*$, i.e.\ with global gauge group
$(SU(3)\times SU(2)\times U(1)_Y)/\mathbb Z_6$ --- developing the
article's recorded open pointer into an explicit congruence.
\emph{Containment of the recorded refinement.}
At $t=0$ (colour singlets) C-$\Gamma$ reduces exactly to the
recorded electroweak $\mathbb Z_2$ condition $2T+Y\in2\mathbb Z$;
C-$\Gamma$ is its colour-completed extension.
\emph{Electric-class corollary.}
With $Q=T_3+Y/2$, C-$\Gamma$ gives $Q\equiv t/3\pmod 1$: colour
singlets carry integer electric charge, (anti)triplets carry
$Q\equiv\mp\tfrac13\ (\mathrm{mod}\ 1)$ --- the
interlock \emph{rule} correlating allowed electric classes with
colour representations.
If the observed quark charges are admitted as input, the minimal
electric lattice compatible with them has unit $e/3$; the
nontrivial content is the correlation rule, not the arithmetic
unit.
\emph{Hedges.}
(i)~Consistency with $\Gamma=\mathbb Z_6$ does not select
$\Gamma$: the observed local spectrum is equally consistent with
the $\mathbb Z_3$, $\mathbb Z_2$, or trivial quotients;
distinguishing them requires extended objects/line-operator data
(the article's recorded citation pointer).
(ii)~For the OP-GUT1 colour completion this is a requirement: the
emergent colour module must realise the $\mathbb Z_3$ centre with
exactly this triality--charge correlation.

\subsection*{2. L2 computed: lattice/interlock action on the
$B{-}L$ flat direction, by branches}
Within the recorded anomaly solution structure of the
one-generation system:
\emph{Branch (i): no $\nu^c$.}
The recorded anomaly system fixes the ratios; the residual
freedom is the overall rescaling $Y\to\lambda Y$.
Imposing C-$\Gamma$ on all five multiplets discretises it:
from $Q$, $\tfrac23+1+\tfrac\lambda3\in2\mathbb Z\Rightarrow
\lambda\equiv1\pmod 6$; the conditions from $u^c,d^c,L,e^c$ are
then satisfied identically.
Residual family: $\lambda\in1+6\mathbb Z$; the observed
assignment is the minimal representative $\lambda=1$, and
\emph{minimality is a convention-level selection}, exactly as in
the recorded OP-GUT2 minimality convention.
\emph{Branch (ii): $Y_{\nu^c}=0$.}
The neutral singlet satisfies C-$\Gamma$ trivially; the analysis
and residual family coincide with branch~(i).
\emph{Branch (iii): free-hypercharge $\nu^c$ (where, and only
where, the recorded $B{-}L$ flat direction appears).}
The recorded flat direction is $Y\to Y+c\,(B{-}L)$ with
$(B{-}L)=(\tfrac13,-\tfrac13,-\tfrac13,-1,1,1)$ for
$(Q,u^c,d^c,L,e^c,\nu^c)$.
Given a base assignment satisfying C-$\Gamma$, the shifted one
satisfies it iff $c\,(B{-}L)_f\in2\mathbb Z$ for every multiplet:
the lepton entries force $c\in2\mathbb Z$, the quark entries
force $c\in6\mathbb Z$.
\emph{Result: the interlock discretises the flat direction to
$c\in6\mathbb Z$ but does not eliminate it.}
Combined residual family in branch (iii):
$(\lambda,c)\in(1+6\mathbb Z)\times6\mathbb Z$, a discrete
two-parameter lattice with minimal representative
$(\lambda,c)=(1,0)$ = the observed assignment; again the
selection of the minimal representative is convention-level.

\subsection*{3. Consequence for the OP-GUT3b statement}
The computations sharpen OP-GUT3b from ``derive the observed
charge assignments'' to: \emph{select the observed point within
an explicitly characterised discrete residual family}
(branch-dependent: $\lambda\in1+6\mathbb Z$, plus $c\in6\mathbb Z$
in the free-$\nu^c$ branch), where anomalies fix the rational
structure, the lattice/interlock fixes the congruence classes,
and the residual selection must come from completion dynamics
(OP-GUT1 module structure), a charge-minimality principle
(convention-level), or structure beyond the present class.
This is a constraint-and-localisation result of the same type as
the OP3.1a column: it narrows where the answer can live without
claiming it.

\subsection*{4. What is NOT claimed}
No derivation of $e_0=e/3$ (observed classes are input); no
selection of $\Gamma$; no derivation of the observed
$(\lambda,c)=(1,0)$; no statement about $g'$, $Z_B^{(0)}$, $k_Y$,
or $\sin^2\theta_W$ (N3, owned by OP-GUT4/OP-GUT7); no insertion;
no status changes.

\subsection*{5. Risk register}
Forbidden: ``ECT derives $e/3$ / the hypercharges / the global
gauge group''; presenting minimality as dynamics; importing
$\Gamma=\mathbb Z_6$ as established; any N3 leakage into
OP-GUT3b text.
Safe: ``consistency requirement on completions, not a
derivation''; ``selection within an explicitly characterised
discrete family''; ``minimality is convention-level''.

\subsection*{6. Questions for GPT}
(Q1)~Verify the congruence algebra independently: the C-$\Gamma$
form $3Y+3d+2t\equiv0\ (\mathrm{mod}\ 6)$, the six verification
values, the $t=0$ containment of the recorded $2T+Y\in2\mathbb Z$,
the electric-class corollary $Q\equiv t/3\ (\mathrm{mod}\ 1)$,
and the two mod-arithmetic results $\lambda\in1+6\mathbb Z$ and
$c\in6\mathbb Z$.
(Q2)~Approve the branch-(iii)-only placement of the $B{-}L$
discussion (per your branch-separation requirement)?
(Q3)~Eventual insertion, after your sign-off: the agreed compact
block ``Hypercharge normalisation: distinct notions and lattice
interlock'' next to the anomaly solution-structure /
OP-GUT3a/3b discussion, containing N1--N4 in two sentences,
C-$\Gamma$ with the verification line, the $t=0$ containment,
the electric-class rule, the branch-resolved residual families
($\lambda$, $c$), the mandatory framing sentence, and a one-line
OP-GUT3a/3b registry echo --- acceptable scope?
(Q4)~The article already cites the global-structure reference in
the recorded pointer; for the new block, is reusing that recorded
citation sufficient, or do you want a page-level verification
round first?
